\chapter{Corpo Docente e de Mediação Pedagógica}

\section{Perfil do corpo docente}

O corpo docente do curso de CDIA é composto por profissionais com sólida formação acadêmica, majoritariamente mestres e doutores na área de Computação e campos correlatos, todos trabalhando por mais de 10 anos em regime de Dedicação Exclusiva na UTFPR (Quadros \ref{quadro3.1} e \ref{quadro3.2}). A qualificação e adequação do corpo docente ao curso é evidenciada por uma produção científica robusta e atualizada, com interação em comunidades científicas internacionais, tendo resultado em publicações em periódicos e conferências de alto impacto voltadas à Ciência de Dados, Aprendizado de Máquina e IA (Quadro \ref{quadro3.3}). Na área aplicada da Computação, os docentes trazem uma vasta experiência profissional aplicada em diferentes áreas da Computação (Quadro \ref{quadro3.4}). Essa combinação entre o rigor acadêmico e a prática de mercado permite que o ensino ultrapasse a teoria, oferecendo aos estudantes uma visão pragmática das demandas do setor tecnológico atual.

No âmbito pedagógico, o corpo docente demonstra ampla proficiência no magistério superior e na modalidade de ensino a distância (EaD), apresentando total domínio das tecnologias educacionais e plataformas digitais de aprendizagem (Moodle e Plataforma Google For Education). Tendo todos os docentes atuado na UTFPR por mais de 10 anos (Quadro \ref{quadro3.1}), os professores utilizam o ambiente virtual de forma estratégica para promover a interatividade e o acompanhamento tempestivo dos discentes, adaptando metodologias ativas às particularidades das unidades curriculares remotas. Essa expertise no ensino mediado por tecnologia garante que a transposição didática de temas complexos da Computação ocorra de forma fluida, garantindo que o estudante seja o protagonista de sua jornada de aprendizagem enquanto conta com a supervisão de especialistas que dominam as ferramentas de suporte pedagógico digital.

Complementando o perfil docente, observa-se uma atuação engajada em projetos de extensão e iniciativas de inovação tecnológica que conectam a instituição à comunidade e ao ecossistema produtivo. Os docentes lideram frentes de pesquisa aplicada e desenvolvimento, transformando o conhecimento gerado em sala de aula em soluções tecnológicas reais e patentes, fomentando o empreendedorismo e a transferência de tecnologia. Esses projetos de extensão não apenas enriquecem o currículo do curso, mas permitem que os professores atuem como mentores em cenários reais de resolução de problemas, garantindo que o corpo docente permaneça na vanguarda das tendências da Inteligência Artificial e da Ciência de Dados, retroalimentando o processo de ensino com experiências práticas de vanguarda. O Quadro \ref{quadro3.5} mostra os projetos de extensão do corpo docente.

Considerando o perfil dos docentes do curso e as necessidades de contratação de docentes específicos das áreas de Ciência de Dados e Inteligência Artificial, também são previstas as seguintes contratações de docentes por área do curso (Quadro \ref{quadro3.6}).

\begin{quadro}[ht]
\caption{Formação, titulação, regime de trabalho e tempo de experiência no magistério superior}
\label{quadro3.1}
\centering

\begin{tabular}{|p{3cm}|p{3cm}|p{3cm}|c|p{2cm}|}
\hline
\textbf{Docente} & \textbf{Graduação} & \textbf{Titulação} & \textbf{Regime} & \textbf{Tempo na IES} \\ \hline
Adriana Herden & Processamento de Dados & Doutorado em Informática Aplicada & DE & 13 anos \\ \hline
Alessandro Botelho Bovo & Ciência da Computação & Doutorado em Eng. e Gestão do Conhecimento & DE & 13 anos \\ \hline
Cláudia Santos Fiuza Lima & Engenharia Elétrica & Doutorado em Engenharia Mecânica & DE & 13 anos \\ \hline
Danilo Sipoli Sanches & Ciência da Computação & Doutorado em Engenharia & DE & 13 anos \\ \hline
Eidy Leandro Tanaka Guandeline & Ciência da Computação & Mestrado em Ciência da Computação & DE & 13 anos \\ \hline
Elias Canhadas Genvigir & Processamento de Dados & Doutorado em Computação Aplicada & DE & 13 anos \\ \hline
Fernando Henrique Campos & Informática & Doutorado em Agronomia & DE & 12 anos \\ \hline
Henrique Yoshikazu Shishido & Tec. em Análise e Desenv. de Sistemas & Doutorado em Ciência da Computação & DE & 13 anos \\ \hline
Marcos Massaki Imamura & Engenharia Elétrica e Eletrônica & Doutorado em Engenharia de Produção & DE & 13 anos \\ \hline
Maurício Correia Lemes Neto & Processamento de Dados & Mestrado em Ciência da Computação & DE & 13 anos \\ \hline
Willian Massami Watanabe & Ciência da Computação & Doutorado em Ciência da Computação & DE & 13 anos \\ \hline
\end{tabular}
\small \\ Fonte: Autoria própria (2026)
\end{quadro}

\begin{quadro}[ht]
\caption{Porcentagem de professores de acordo com sua titulação máxima}
\label{quadro3.2}
\centering
\begin{tabular}{|l|c|c|}
\hline
\textbf{Titulação} & \textbf{Quantidade} & \textbf{\%} \\ \hline
Doutores & 8 & 80\% \\ \hline
Mestres & 2 & 20\% \\ \hline
Especialistas & 0 & 0\% \\ \hline
\textbf{Total} & \textbf{10} & \textbf{100\%} \\ \hline
\end{tabular}
\small \\ Fonte: Autoria própria (2026)
\end{quadro}

\begin{quadro}[ht]
\caption{Experiência internacional e Produção Intelectual/Técnica dos docentes}
\label{quadro3.3}
\centering
\begin{tabular}{|p{4cm}|p{6cm}|c|}
\hline
\textbf{Docente} & \textbf{Experiência internacional} & \textbf{Produção (3 anos)} \\ \hline
Adriana Herden & Artigos em eventos e periódicos internacionais & 6 \\ \hline
Alessandro Botelho Bovo & Pesquisador na Leiden University (Holanda) & 3 \\ \hline
Danilo Sipoli Sanches & Pós-Doutorado na Colorado State University (EUA) & 20 \\ \hline
Elias Canhadas Genvigir & Professor Visitante na Universitat de València & 0 \\ \hline
Henrique Yoshikazu Shishido & Artigos em eventos e periódicos internacionais & 4 \\ \hline
Marcos Massaki Imamura & Artigos em eventos e periódicos internacionais & 5 \\ \hline
Willian Massami Watanabe & Artigos em eventos e periódicos internacionais & 9 \\ \hline
\end{tabular}
\small \\ Fonte: Autoria própria (2026)
\end{quadro}

\begin{quadro}[ht]
\caption{Atuação profissional na área aplicada do curso por docente}
\label{quadro3.4}
\centering
\begin{tabular}{|p{5cm}|p{4cm}|c|}
\hline
\textbf{Docente} & \textbf{Experiência aplicada} & \textbf{Tempo de vínculo} \\ \hline
Adriana Herden & Analista de Sistemas & 3 anos \\ \hline
Danilo Sipoli Sanches & Engenheiro de Software & 2 anos \\ \hline
Eidy Leandro Tanaka Guandeline & Analista de Sistemas & 6 anos \\ \hline
Elias Canhadas Genvigir & Engenheiro de Software & 2 anos \\ \hline
Henrique Yoshikazu Shishido & Engenheiro de Software & 3 anos \\ \hline
Maurício Correia Lemes Neto & Engenheiro de Software & 6 anos \\ \hline
Willian Massami Watanabe & Engenheiro de Software & 2 anos \\ \hline
\end{tabular}
\small \\ Fonte: Autoria própria (2026)
\end{quadro}

\begin{quadro}[ht]
\caption{Projetos de extensão e inovação tecnológica dos docentes}
\label{quadro3.5}
\centering
\begin{tabular}{|p{4cm}|p{10cm}|}
\hline
\textbf{Docente} & \textbf{Projeto de extensão e inovação tecnológica} \\ \hline
Danilo Sipoli Sanches & Cooperação Técnica HUB IA SENAI/PR; Teoria e Prática no desenv. da capacidade empreendedora \\ \hline
Fernando Henrique Campos & Inclusão do Jogo de Xadrez na UTFPR; Comparação da intuitividade em sistemas móveis \\ \hline
Marcos Massaki Imamura & Protótipo off Road 4x4: Equipe Londribaja \\ \hline
Willian Massami Watanabe & Acessibilidade para Deficientes; Projetos de Código Aberto (Reactstrap) \\ \hline
\end{tabular}
\small \\ Fonte: Autoria própria (2026)
\end{quadro}

\begin{quadro}[ht]
\caption{Previsão de contratação docente por área do curso}
\label{quadro3.6}
\centering
\begin{tabular}{|p{8cm}|c|}
\hline
\textbf{Perfil e área de atuação do docente} & \textbf{Quantidade} \\ \hline
Visão Computacional & 2 \\ \hline
Processamento de Linguagem Natural & 2 \\ \hline
Aprendizado de Máquina & 4 \\ \hline
Inteligência Artificial & 3 \\ \hline
Big Data e Computação em Nuvem & 2 \\ \hline
Matemática & 1 \\ \hline
Estatística & 1 \\ \hline
\textbf{TOTAL} & \textbf{15} \\ \hline
\end{tabular}
\small \\ Fonte: Autoria própria (2026)
\end{quadro}

\section{Corpo de Mediação Pedagógica}

O curso de CDIA estrutura-se com unidades curriculares com carga horária na modalidade a distância (EaD) e estabelece que as turmas devem conter, no máximo, 40 alunos, garantindo que o próprio professor atue diretamente na mediação pedagógica. Essa configuração é estratégica para favorecer a integração plena do mediador com os conteúdos lecionados e promover uma interatividade superior entre docentes e discentes. Alinhado à legislação vigente, este corpo de mediação possui formação e expertise no mesmo campo de conhecimento do curso, além de expertise com as tecnologias digitais de mediação pedagógica (Moodle e \textit{Google for Education Suite}) conforme destacado no Perfil do Corpo Docente. Essa característica de docentes mediadores permite interações com a tempestividade, profundidade e abrangência inerentes ao processo educacional, respeitando o quantitativo de estudantes e as metodologias propostas, que só seriam possíveis nessa configuração de curso.

As disciplinas com carga horária EaD do curso apresentam uma característica prática inerente, sendo essa carga horária dedicada, muitas vezes, para implementação e desenvolvimento de projetos. Nesse cenário, o docente-mediador planeja e executa as atividades, estabelecendo espaços de interação, apresentando exemplos práticos relacionados à realidade local dos grupos de estudantes e associados aos componentes curriculares. Os projetos promovem o engajamento em atividades colaborativas e incentiva o protagonismo estudantil, sendo o mediador constantemente avaliado pelos alunos quanto à qualidade e rapidez de seu suporte. A experiência profissional e acadêmica do corpo de mediação assegura que o processo de ensino-aprendizagem seja tanto tecnicamente rigoroso quanto pedagogicamente acessível.


\section{Ambientes de aprendizagem em educação profissional e tecnológica}

Os ambientes de aprendizagem do curso são projetados para garantir a transposição efetiva entre teoria e prática, característica essencial na área de Ciência da Computação. O curso dispõe de laboratórios modernos com ampla disponibilidade de computadores de alto desempenho, onde o Núcleo Docente Estruturante (NDE) realiza análises e avaliações periódicas para assegurar a atualização permanente do parque tecnológico e o alinhamento das práticas laborais às inovações do setor. Essa infraestrutura permite que os estudantes desenvolvam habilidades e competências alinhadas ao perfil do egresso. Esse monitoramento constante pelo NDE garante que os laboratórios não sejam apenas espaços físicos, mas ecossistemas de inovação que mimetizam os desafios reais do mercado de trabalho.

A integração com o setor produtivo é evidenciada pelos projetos de extensão e inovação tecnológica liderados pelos docentes, que frequentemente envolvem cooperação direta com empresas de tecnologia. O curso mantém parcerias estratégicas com o consolidado arranjo produtivo de TI de Londrina, aproveitando a densidade tecnológica da região para o engajamento dos alunos. A presença de empresas de tecnologia na região, como Tata Consultancy Services (TCS), Atos, Trimble, ID Tech e Unico ID, cria um ambiente propício para estágios, projetos conjuntos e vivência profissional. Essas parcerias asseguram que o currículo seja retroalimentado pelas demandas de empresas líderes, permitindo que as atividades de extensão funcionem como pontes diretas para o desenvolvimento de soluções de software e inteligência artificial aplicadas ao mercado real.

Na UTFPR, o desempenho docente é orientado por práticas que asseguram coerência pedagógica, atualização contínua e compromisso com a formação profissional tecnológica. Cada docente atua de acordo com as atribuições descritas no PPC, participando ativamente da análise e revisão dos objetivos, dos conteúdos e das atividades de aprendizagem das Unidades Curriculares (UCs). Além disso, apresenta aos estudantes, presencialmente ou por meio de tecnologias digitais, o plano de ensino das UCs que ministra, atualizando e aprimorando esse plano e seus planos de aula com base nos resultados de sua autoavaliação.

O corpo docente utiliza recursos pedagógicos, tecnológicos, digitais e laboratoriais alinhados às necessidades de aprendizagem do curso, adotando metodologias ativas que favorecem o protagonismo estudantil e o desenvolvimento de competências,  com ênfase no uso de Aprendizado Baseado em Projetos (PBL - Project-Based Learning). As UCs são permanentemente articuladas a publicações e referenciais atualizados, garantindo uma formação conectada às transformações da área de Ciência de Dados e Inteligência Artificial. Os docentes também promovem atividades formativas que fortalecem o perfil profissional esperado, contribuindo para que os egressos atuem de maneira crítica, ética, inovadora e tecnicamente qualificada diante dos desafios contemporâneos da IA.

A cultura de avaliação é um pilar central para a melhoria contínua do curso. No âmbito de avaliação interna, o corpo docente é sistematicamente avaliado pelos estudantes em relação ao seu desempenho didático-pedagógico, além de realizar processos de autoavaliação para o aprimoramento constante de sua prática. O processo de avaliação de desempenho docente considera, de forma integrada, o desempenho nas unidades curriculares (considerando a evasão/retenção) e desenvolvimento de competências dos estudantes, servindo de base para ajustes curriculares e pedagógicos. De forma complementar, as avaliações externas do SINAES, incluindo os processos de Reconhecimento de Curso e o ENADE, asseguram que o curso mantenha padrões de excelência reconhecidos pelo Ministério da Educação.

A produção do corpo docente é rigorosamente alinhada aos objetivos do curso, sendo comprovada por publicações em meios de difusão de alto impacto e relevância científica na área de CDIA. Tais produções não se limitam à teoria, gerando frequentemente propriedade intelectual e direitos autorais que reforçam o caráter inovador da instituição. Além disso, as atividades de pesquisa e produção técnica envolvem ativamente a participação de estudantes e promovem ações diretas com a comunidade por meio de atividades de extensão. Essa articulação garante que o conhecimento produzido pelos docentes transite entre a academia e a sociedade, consolidando um perfil docente que é, simultaneamente, um pesquisador de ponta e um mentor engajado com a realidade social e tecnológica.

\section{Formação e Desenvolvimento Permanente do Corpo Docente}

A formação continuada dos professores da UTFPR é estruturada pelo Programa de Desenvolvimento Profissional Docente (PDPD), regulamentado pelas Resoluções COGEP nº 32/2019 e nº 44/2020 \cite{UTFPR_Desenvolvimento_2019}. Esse programa constitui um plano institucional integrado de formação para a docência, com a finalidade de qualificar permanentemente as práticas de ensino e fortalecer a identidade docente no contexto da educação tecnológica. O PDPD está organizado em dois planos: o Desenvolvimento Profissional Docente Inicial (PD$^2$i), destinado aos professores ingressantes e em estágio probatório, e o Desenvolvimento Profissional Docente Continuado (PD$^2$c), voltado aos docentes estáveis, assegurando a atualização sistemática em temas pedagógicos, metodológicos, tecnológicos e profissionais.

O programa contempla módulos de formação que abordam a universidade e o trabalho docente, princípios da educação tecnológica, demandas contemporâneas do ensino superior, relações profissionais, metodologias e didática, tecnologias educacionais e avaliação da aprendizagem. Além desses módulos, a formação pode incluir participação em eventos acadêmicos, grupos de estudo, ações de pesquisa e extensão, projetos pedagógicos inovadores e acompanhamento pedagógico especializado. Na UTFPR, o PDPD assume papel estratégico ao apoiar o docente na incorporação de tecnologias digitais, metodologias ativas, práticas de ensino baseadas em dados e inovação acadêmica, garantindo que a atuação docente acompanhe as transformações científicas, tecnológicas e sociais que caracterizam a formação em ciência, tecnologia, engenharia, matemática e inteligência artificial. Dessa forma, a UTFPR assegura um processo contínuo de aperfeiçoamento que impacta diretamente a qualidade acadêmica e o alinhamento pedagógico dos cursos.

A PROGRAD oferece uma série de eventos para promover a formação continuada do docente. Alguns dos eventos institucionais de maiores destaques são os fóruns de coordenadores, as oficinas de design de cursos e as oficinas de design de unidades curriculares.

Nos fóruns de coordenadores, reúne-se docentes de todos os campi para discutir temas como evasão, retenção, mobilidade acadêmica, incorporação de atividades de extensão, EaD, legislação vigente e, principalmente, modelos de inovação curricular e diferentes metodologias de ensino e aprendizagem. As discussões têm como objetivo dar suporte para a implementação de melhorias nos cursos de graduação da UTFPR.

As oficinas de design de cursos e de unidades curriculares são uma série de encontros de formação e elaboração cooperativa para os membros dos NDEs, afim de promover a reformulação de seus respectivos cursos e unidades curriculares, por meio de uma metodologia baseada em competências. Esta vivência propicia a reflexão sobre as práticas de ensino e monitoramento de índices como os de aprovação e de evasão. Além disso, as metodologias de ensino e avaliação juntamente com os instrumentos pedagógicos utilizados em aula são revistos e aperfeiçoados a partir, principalmente, do trabalho colaborativo e do compartilhamento de informações da rede que se formou.

A UTFPR também oferece mecanismos para o desenvolvimento profissional do corpo docente na forma de licenças para pós-graduação; licenças capacitação; e atividades de aperfeiçoamento técnico e pedagógico. Em relação às licenças para pós-graduação stricto sensu, a Instrução Normativa Conjunta PROPPG/DIRGET nº 01/2019, de 03 de outubro de 2019, determina que o afastamento para a realização de pós-graduação stricto sensu ocorra por meio da aprovação em edital de seleção publicado de acordo com a política institucional de capacitação expressa no Plano de Desenvolvimento de Pessoas (PDP). A licença capacitação possibilita ao servidor participar de ações de desenvolvimento técnico, comportamental, científico ou cultural por até 3 (três) meses a cada quinquênio de efetivo exercício. As normas para sua concessão estão estabelecidas no regulamento aprovado em Deliberação do COUNI nº  32, de 20/12/2019. Finalmente, as atividades de aperfeiçoamento técnico e pedagógico ocorrem por meio da participação dos docentes em eventos, seminários, simpósios e fóruns científicos nacionais e internacionais. Existe, ainda, a semana de planejamento docente, processo que, normalmente, ocorre semestralmente, e é oferecido pela própria instituição. Neste período, são oferecidos palestras e mini-cursos, ministrados por convidados externos e docentes do próprio campus.


\section{Formação e Desenvolvimento Permanente do Corpo de Mediação Pedagógica}

Considerando a estrutura pedagógica do curso de CDIA, em que o próprio corpo docente assume a responsabilidade direta pela mediação das unidades curriculares com carga horária a distância (EaD), as políticas de qualificação da UTFPR são aplicadas de forma integrada. Dessa maneira, as ações de formação continuada promovidas pelo Programa de Desenvolvimento Profissional Docente (PDPD), incluindo os planos PD$^2$i e PD$^2$c, licenças para pós-graduação, licenças capacitação e atividades de aperfeiçoamento técnico e pedagógico são também incentivadas aos mediadores, garantindo que o suporte pedagógico no ambiente virtual de aprendizagem possua o mesmo rigor técnico e metodológico do ensino presencial. Essa equivalência assegura que os docentes-mediadores participem dos fóruns de coordenadores, das oficinas de design de unidades curriculares e das semanas de planejamento, capacitando-os para o uso estratégico de tecnologias educacionais e para a aplicação de práticas de ensino baseadas em dados, o que consolida a indissociabilidade entre a regência da disciplina e a mediação pedagógica no contexto da educação tecnológica.

A organização acadêmica e administrativa do curso de CDIA observa integralmente as disposições da Resolução nº 145/2019 – COGEP, que regulamenta o processo de escolha, designação, requisitos e atribuições dos Coordenadores de Curso de Graduação na Universidade. Conforme essa normativa, a coordenação deve ser exercida por docente pertencente ao quadro efetivo, com formação e experiência compatíveis, atendendo aos requisitos institucionais de elegibilidade, dedicação e atuação prévia no curso.

A coordenação do curso de CDIA é exercida por um docente com regime de trabalho de Dedicação Exclusiva, garantindo carga horária dedicada e adequada para a gestão plena do curso, a participação ativa em órgãos colegiados e o atendimento a docentes, estudantes e mediadores pedagógicos. A atuação do coordenador é pautada no Projeto Pedagógico do Curso (PPC) e executada por meio de um plano de ação documentado, compartilhado com a comunidade acadêmica. Esse plano é de conhecimento público, alcançando estudantes e profissionais em todos os ambientes de aprendizagem, e sua execução é monitorada por indicadores de desempenho (como taxas de evasão, retenção e efetividade, inserção profissional dos alunos, atratividade do curso, entre outros), disponíveis para consulta pela comunidade.

No exercício de suas funções, o coordenador faz uso estratégico de tecnologias educacionais digitais para assegurar uma interação fluida com gestores, docentes e alunos, especialmente considerando a existência de unidades curriculares EaD. Sua gestão foca na potencialização do corpo docente e dos mediadores, buscando uma alocação que respeite o perfil de cada profissional em sintonia com as competências desejadas para o egresso. Além de favorecer a integração contínua entre todos os atores do processo educativo, o coordenador atua como um catalisador da melhoria contínua dos processos de ensino e aprendizagem, promovendo inovações pedagógicas e curriculares que mantenham o curso na vanguarda tecnológica.

A instituição viabiliza o suporte necessário para a excelência dessa função por meio de programas de formação e atualização permanente, capacitando o coordenador para a utilização de novos recursos digitais, serviços e equipamentos. Esse processo de formação continuada é sistematicamente avaliado pela Comissão Própria de Avaliação (CPA), cujos resultados geram subsídios para ações de aprimoramento da gestão. Dessa forma, a coordenação não apenas atende às demandas administrativas, mas lidera o curso com uma visão técnica e pedagógica alinhada aos desafios da Inteligência Artificial, assegurando o cumprimento das metas institucionais e a qualidade acadêmica reconhecida pelo SINAES.

Dessa forma, o curso se organiza em consonância com as diretrizes institucionais vigentes, assegurando gestão acadêmica qualificada, continuidade administrativa e alinhamento ao PDI da UTFPR.


\section{Núcleo Docente Estruturante}

O Núcleo Docente Estruturante (NDE) do curso, é um órgão consultivo da coordenação responsável pela concepção, consolidação, acompanhamento e permanente atualização do Projeto Pedagógico do Curso, assegurando sua coerência acadêmica, pedagógica e institucional. Sua composição e funcionamento seguem as determinações da Resolução interna nº 009/12-COGEP, de 13 de abril de 2012 \cite{UTFPR_NDE_2012}. No âmbito deste curso tecnológico voltado às áreas STEM e à Inteligência Artificial, o NDE desempenha papel estratégico ao garantir que a formação permaneça atualizada frente às inovações científicas e tecnológicas, às demandas do setor produtivo e às diretrizes institucionais da UTFPR.

O Núcleo Docente Estruturante (NDE) do curso é constituído por um colegiado de excelência acadêmica, composto por, no mínimo, cinco docentes, incluindo o coordenador do curso. Em estrita observância à legislação vigente \cite{CONAES_NDE_2010}, sendo que, no mínimo, 20\% do grupo possui dedicação integral e pelo menos 60\% detém titulação stricto sensu. Para garantir a memória institucional e a estabilidade das diretrizes pedagógicas, o NDE mantém a continuidade de parte de seus membros entre ciclos avaliativos, assegurando que o acompanhamento do curso seja contínuo e que as Diretrizes Curriculares Nacionais (DCN) \cite{MEC_2016}, quando aplicáveis, sejam integralmente cumpridas e consolidadas no cotidiano acadêmico.

A atuação do NDE é pautada na revisão periódica e atualização do Projeto Pedagógico do Curso (PPC), fundamentada em estudos que validam a adequação do perfil do egresso tanto às normas educacionais quanto às dinâmicas exigências do mundo do trabalho em tecnologia. Por meio de ferramentas e recursos digitais avançados, o núcleo promove a interatividade com o corpo docente e mediadores pedagógicos, zelando pela interdisciplinaridade e pelo constante aperfeiçoamento didático-pedagógico. O grupo atua proativamente no incentivo a atividades de ensino, pesquisa e extensão, garantindo que os conteúdos e metodologias adotados possuam aderência real às necessidades profissionais de um mercado em constante transformação pela Inteligência Artificial.

Comprometido com a cultura da qualidade, o NDE promove a autoavaliação sistemática do curso, utilizando instrumentos próprios e os relatórios da Comissão Própria de Avaliação (CPA) \cite{UTFPR_CPA_2004, UTFPR_CPA_2009} para subsidiar a atualização contínua do PPC e das metodologias de ensino. Os resultados das avaliações de aprendizagem são rigorosamente analisados para ajustar as estratégias pedagógicas e garantir o desenvolvimento das competências previstas. Além disso, os membros do NDE mantêm-se em constante processo de atualização profissional através de interações diretas com egressos, empreendedores e especialistas externos atuantes no setor de TI, assegurando que a formação oferecida pela instituição esteja plenamente alinhada às tendências de inovação e práticas de mercado mais recentes.

\section{Colegiado do Curso}

O Colegiado do curso é instituído conforme a Resolução COGEP nº 103/2019, constitui-se como órgão propositivo responsável por assessorar a Coordenação em assuntos relacionados às políticas de ensino, pesquisa e extensão, assegurando coerência com os princípios, finalidades e objetivos institucionais da UTFPR. Sua organização, composição e funcionamento abrangem atribuições como elaboração e atualização do Projeto Pedagógico do Curso, acompanhamento da implementação curricular, emissão de pareceres acadêmicos, definição de critérios para atividades extensionistas, estágios, TCC, atividades complementares e mobilidade acadêmica, bem como apoio aos processos de regulação, avaliação e desenvolvimento institucional. O colegiado é composto pela Coordenação e pelos responsáveis por estágio, TCC, extensão, internacionalização e atividades complementares, além de representantes docentes e discentes, observando os critérios eletivos e mandatários definidos na Resolução. Reúne-se periodicamente, delibera por maioria simples e mantém registro formal de suas atividades, garantindo gestão colegiada, transparência acadêmica e aderência plena às normativas institucionais \cite{UTFPR_Colegiado_2019}.

O colegiado do curso de CDIA é o órgão consultivo e deliberativo devidamente institucionalizado, regido por regimento próprio e com planejamento sistemático de suas atividades. Sua composição garante a representatividade democrática de docentes das diversas unidades curriculares, estudantes e dos próprios professores que atuam na mediação pedagógica, estimulando práticas interdisciplinares e a integração entre diferentes eixos do conhecimento. Com reuniões periódicas registradas em atas e apoiadas por um sistema de suporte ao registro e execução de decisões, o Colegiado zela pelo planejamento acadêmico nos âmbitos do ensino, pesquisa e extensão. Além de acompanhar a implementação de práticas de gestão e a atuação do NDE e da coordenação, o órgão demonstra uma atuação estratégica ao estabelecer políticas de melhoria contínua, sempre pautando suas deliberações no desenvolvimento das competências e no perfil do egresso, garantindo que a formação técnica e humana esteja em plena harmonia com as constantes transformações do mundo do trabalho e as demandas da sociedade contemporânea.